% GENERATED FILE. Edit canonical lessons, then run npm run generate:books.
% canonical-source-hash: fnv1a64:96456dbbf016d482
% canonical-lessons: FR-C349-provenir, FR-C349-resulter, FR-C349-influence, FR-C349-influencer, FR-C349-engendrer, FR-R349-revision-342-349

\chapter{D'où ça vient}
\label{ch:fr-d-ou-ca-vient}
\hlchaptermodality{349}

\begin{chapteropening}
\textbf{By the end of this chapter:} \emph{I can say where something comes from, what it results from, what influences it and what it gives rise to.}

Complete the last lesson of chapter 349: I can say where something comes from, what it results from, what influences it and what it gives rise to.
\end{chapteropening}

\section[provenir]{provenir — to come from, to originate in}

\label{lesson:FR-C349-provenir}



\begin{quote}
Before the new one: say the French for from then on, then the French for the origin.
\end{quote}

\subsection*{You'll want to know: provenir}
\textbf{provenir} — \textquotedblleft{}to come from, to originate in\textquotedblright{}.

\begin{quote}
Ce bruit provient de la cuisine.
\end{quote}

\begin{quote}
\emph{This noise is coming from the kitchen.}
\end{quote}

Conjugated like \emph{venir}, with \emph{être}. \emph{En provenance de} — arriving from.

\begin{cousinweb}
\emph{provenir} comes from Latin \emph{prōvenīre} \textquotedblleft{}to come forth, to arise\textquotedblright{} (EtymDB; Lewis \& Short), built on \emph{venīre} \textquotedblleft{}to come\textquotedblright{}.
\end{cousinweb}

\begin{grammarlens}[title={What you've built}]
One more word for telling, explaining and describing at B1.
\end{grammarlens}

\subsection*{Your turn}
\begin{itemize}
\raggedright
  \item \emph{Say it:} \emph{provenir}
  \item \emph{Say it:} \emph{provenir}, once more
  \item \emph{Say it:} a sentence of your own with \emph{provenir}
  \item \emph{Recall:} say the French for from then on, then the French for the origin, then say \emph{provenir} again
\end{itemize}

\begin{tcolorbox}[breakable,colback=teal!4,colframe=teal!35!black,title={Before you move on}]
What does provenir mean? (\textquotedblleft{}To come from, to originate in\textquotedblright{}.) Say it once more.
\end{tcolorbox}
\section[résulter]{résulter — to result (from)}

\label{lesson:FR-C349-resulter}



\begin{quote}
Before the new one: say the French for the origin, then the French for to come from.
\end{quote}

\subsection*{You'll want to know: résulter}
\textbf{résulter} — \textquotedblleft{}to result (from)\textquotedblright{}.

\begin{quote}
Cet accident résulte d'une erreur humaine.
\end{quote}

\begin{quote}
\emph{This accident results from human error.}
\end{quote}

\emph{Il en résulte que} — the result is that. Mostly in the third person.

\begin{cousinweb}
\emph{résulter} is a learned borrowing from Latin \emph{resultāre} \textquotedblleft{}to spring back, to rebound\textquotedblright{} (Lewis \& Short), built on \emph{saltāre} \textquotedblleft{}to leap\textquotedblright{}.
\end{cousinweb}

\begin{grammarlens}[title={What you've built}]
One more word for telling, explaining and describing at B1.
\end{grammarlens}

\subsection*{Your turn}
\begin{itemize}
\raggedright
  \item \emph{Say it:} \emph{résulter}
  \item \emph{Say it:} \emph{résulter}, once more
  \item \emph{Say it:} a sentence of your own with \emph{résulter}
  \item \emph{Recall:} say the French for the origin, then the French for to come from, then say \emph{résulter} again
\end{itemize}

\begin{tcolorbox}[breakable,colback=teal!4,colframe=teal!35!black,title={Before you move on}]
What does résulter mean? (\textquotedblleft{}To result (from)\textquotedblright{}.) Say it once more.
\end{tcolorbox}
\section[l'influence]{l'influence — the influence}

\label{lesson:FR-C349-influence}



\begin{quote}
Before the new one: say the French for to come from, then the French for to result.
\end{quote}

\subsection*{You'll want to know: l'influence}
\textbf{l'influence} — \textquotedblleft{}the influence\textquotedblright{}.

\begin{quote}
Ses amis ont une mauvaise influence sur lui.
\end{quote}

\begin{quote}
\emph{His friends are a bad influence on him.}
\end{quote}

\emph{Sous l'influence de} — under the influence of. Feminine.

\begin{cousinweb}
\emph{influence} is usually traced to Medieval Latin \emph{influentia} \textquotedblleft{}a flowing in\textquotedblright{}, built on Latin \emph{influere} \textquotedblleft{}to flow into\textquotedblright{} (Lewis \& Short), from \emph{fluere} \textquotedblleft{}to flow\textquotedblright{}.
\end{cousinweb}

\begin{grammarlens}[title={What you've built}]
One more word for telling, explaining and describing at B1.
\end{grammarlens}

\subsection*{Your turn}
\begin{itemize}
\raggedright
  \item \emph{Say it:} \emph{l'influence}
  \item \emph{Say it:} \emph{l'influence}, once more
  \item \emph{Say it:} a sentence of your own with \emph{l'influence}
  \item \emph{Recall:} say the French for to come from, then the French for to result, then say \emph{l'influence} again
\end{itemize}

\begin{tcolorbox}[breakable,colback=teal!4,colframe=teal!35!black,title={Before you move on}]
What does l'influence mean? (\textquotedblleft{}The influence\textquotedblright{}.) Say it once more.
\end{tcolorbox}
\section[influencer]{influencer — to influence}

\label{lesson:FR-C349-influencer}



\begin{quote}
Before the new one: say the French for to result, then the French for the influence.
\end{quote}

\subsection*{You'll want to know: influencer}
\textbf{influencer} — \textquotedblleft{}to influence\textquotedblright{}.

\begin{quote}
La publicité influence nos choix.
\end{quote}

\begin{quote}
\emph{Advertising influences our choices.}
\end{quote}

\emph{Nous influençons}, with \emph{ç}.

\begin{cousinweb}
\emph{influencer} is built on \emph{influence}, from Medieval Latin \emph{influentia} \textquotedblleft{}a flowing in\textquotedblright{}, on Latin \emph{influere} \textquotedblleft{}to flow into\textquotedblright{}, from \emph{fluere} \textquotedblleft{}to flow\textquotedblright{} (Lewis \& Short).
\end{cousinweb}

\begin{grammarlens}[title={What you've built}]
One more word for telling, explaining and describing at B1.
\end{grammarlens}

\subsection*{Your turn}
\begin{itemize}
\raggedright
  \item \emph{Say it:} \emph{influencer}
  \item \emph{Say it:} \emph{influencer}, once more
  \item \emph{Say it:} a sentence of your own with \emph{influencer}
  \item \emph{Recall:} say the French for to result, then the French for the influence, then say \emph{influencer} again
\end{itemize}

\begin{tcolorbox}[breakable,colback=teal!4,colframe=teal!35!black,title={Before you move on}]
What does influencer mean? (\textquotedblleft{}To influence\textquotedblright{}.) Say it once more.
\end{tcolorbox}
\section[engendrer]{engendrer — to generate, to give rise to}

\label{lesson:FR-C349-engendrer}



\begin{quote}
Before the new one: say the French for the influence, then the French for to influence.
\end{quote}

\subsection*{You'll want to know: engendrer}
\textbf{engendrer} — \textquotedblleft{}to generate, to give rise to\textquotedblright{}.

\begin{quote}
Le bruit engendre la fatigue.
\end{quote}

\begin{quote}
\emph{Noise brings on tiredness.}
\end{quote}

Formal. The older sense is to father a child.

\begin{cousinweb}
\emph{engendrer} is inherited from Latin \emph{ingenerāre} \textquotedblleft{}to implant, to produce\textquotedblright{} (EtymDB; Lewis \& Short), built on \emph{generāre} \textquotedblleft{}to beget\textquotedblright{}, from \emph{genus} \textquotedblleft{}birth, kind\textquotedblright{}.
\end{cousinweb}

\begin{grammarlens}[title={What you've built}]
One more word for telling, explaining and describing at B1.
\end{grammarlens}

\subsection*{Your turn}
\begin{itemize}
\raggedright
  \item \emph{Say it:} \emph{engendrer}
  \item \emph{Say it:} \emph{engendrer}, once more
  \item \emph{Say it:} a sentence of your own with \emph{engendrer}
  \item \emph{Recall:} say the French for the influence, then the French for to influence, then say \emph{engendrer} again
\end{itemize}

\begin{tcolorbox}[breakable,colback=teal!4,colframe=teal!35!black,title={Before you move on}]
What does engendrer mean? (\textquotedblleft{}To generate, to give rise to\textquotedblright{}.) Say it once more.
\end{tcolorbox}
\section[(dialogue)]{Retour sur les chapitres 342–349}

\label{lesson:FR-R349-revision-342-349}



\begin{quote}
Say the French for to evaluate and to judge.
\end{quote}

\subsection*{Your turn}
Say these aloud:

\begin{itemize}
\raggedright
  \item to evaluate, to judge, the judgement, subjective, impartial
  \item the advantage, the drawback, the aspect, the criticism, to stress
  \item to reason, the reasoning, to conclude, the subtle difference, to qualify
  \item to approve, to disapprove, the supporter, the opponent, to argue a case
  \item to cause, to bring about, the consequence, to lead to, to follow
  \item since, given that, seeing that, consequently, that is why
  \item because of, thanks to, indeed, from then on, the origin
  \item to come from, to result, the influence, to influence, to generate
\end{itemize}

\begin{tcolorbox}[breakable,colback=teal!4,colframe=teal!35!black,title={Before you move on}]
To evaluate? (\textbf{évaluer}.) To judge? (\textbf{juger}.) The judgement? (\textbf{le jugement}.) Subjective? (\textbf{subjectif}.) Impartial? (\textbf{impartial}.) The advantage? (\textbf{l'avantage}.) The drawback? (\textbf{l'inconvénient}.) The aspect? (\textbf{l'aspect}.) The criticism? (\textbf{la critique}.) To stress? (\textbf{souligner}.) To reason? (\textbf{raisonner}.) The reasoning? (\textbf{le raisonnement}.) To conclude? (\textbf{conclure}.) The subtle difference? (\textbf{la nuance}.) To qualify? (\textbf{nuancer}.) To approve? (\textbf{approuver}.) To disapprove? (\textbf{désapprouver}.) The supporter? (\textbf{le partisan}.) The opponent? (\textbf{l'adversaire}.) To argue a case? (\textbf{argumenter}.) To cause? (\textbf{causer}.) To bring about? (\textbf{provoquer}.) The consequence? (\textbf{la conséquence}.) To lead to? (\textbf{entraîner}.) To follow? (\textbf{découler}.) Since? (\textbf{puisque}.) Given that? (\textbf{étant donné que}.) Seeing that? (\textbf{vu que}.) Consequently? (\textbf{par conséquent}.) That is why? (\textbf{c'est pourquoi}.) Because of? (\textbf{à cause de}.) Thanks to? (\textbf{grâce à}.) Indeed? (\textbf{en effet}.) From then on? (\textbf{dès lors}.) The origin? (\textbf{l'origine}.) To come from? (\textbf{provenir}.) To result? (\textbf{résulter}.) The influence? (\textbf{l'influence}.) To influence? (\textbf{influencer}.) To generate? (\textbf{engendrer}.)
\end{tcolorbox}
